% !TeX program = lualatex
% !TeX encoding = utf8
% !TeX spellcheck = uk_UA
% !TeX root = ../EMConspect.tex
% УВАГА: для \enquote{...} потрібен пакет csquotes (\usepackage[ukrainian=quotes]{csquotes}
% в преамбулі EMConspect.tex) --- переконайтесь, що він підключений.

%=========================================================
\chapter{Типи діелектриків та їх застосування}\label{\currfilebase-III}
\graphicspath{{\currfilebase/Pictures}}
%=========================================================
\epigraph{Від ізолятора в кабелі до рідкокристалічного дисплея --- вся техніка навколо нас стоїть на тому, як речовина відгукується на електричне поле.
}{Переосмислення ідей сучасної фізики}

У попередніх двох розділах ми розібрали \emphx{лінійний} діелектрик: його поляризація пропорційна полю, $\vect P = \chi\Efield$, а проникність
$\epsilon = 1 + 4\pi\chi$ --- стала величина, що залежить лише від природи молекул, їхньої концентрації та температури. Проте далеко не всі
діелектрики поводяться так просто. У цьому розділі ми познайомимось із \emphx{особливими} класами діелектриків, чиї властивості широко
використовуються в техніці: електретами, сегнетоелектриками, п'єзо- та піроелектриками, релаксорами, антисегнетоелектриками, рідкими кристалами, а
також пасивними діелектриками-ізоляторами.

%% --------------------------------------------------------
\section{Електрети}
%% --------------------------------------------------------

\emphx{Електрети} --- це діелектрики, здатні тривалий час зберігати поляризацію \emphx{після зняття} зовнішнього поля, яке цю поляризацію створило.
Електрет часто називають \enquote{\emphx{електричним аналогом постійного магніту}}: подібно до того, як постійний магніт створює магнітне поле без
зовнішнього джерела струму, заряджений електрет створює електричне поле $\Efield$ і поле індукції $\Dfield$ поза своєю поверхнею навіть за
відсутності будь-якого зовнішнього поля (рис.~\ref{pic:Electret}).

\begin{figure}[h!]\centering
	\localinput{electret-fields.tikz}
	\caption{Поля зарядженого електрета\label{pic:Electret}}
\end{figure}

\begin{Attention}
	Простежте аналогію між електретом та постійним магнітом (рис.~\ref{pic:Electret}). Поля $\Bfield$ і $\Hfield$ будуть докладно введені в розділі~<<\nameref{MagFieldMedia}>>, а тут ми лише скористаємось ними для порівняння:

	\begin{itemize}
		\item Поле $\Dfield$ (рис.~а) подібне до поля $\Bfield$: лінії обох векторів неперервні й замкнені, без початку й кінця.
		\item Поле $\Efield$ (рис.~б) подібне до поля $\Hfield$: всередині електрета $\Efield$ напрямлений протилежно до $\vect{P}$ (деполяризуюче поле) --- це точний аналог розмагнічувального поля $\Hfield$ всередині магніту.
	\end{itemize}

    \vspace*{1ex}

	В ідеальному електреті вільних зарядів немає \emphx{взагалі}, тож $\rho_{\text{вільні}}\equiv0$ і рівняння $\nabla\cdot\Dfield=\rho_{\text{вільні}}$ перетворюється на тотожність $\nabla\cdot\Dfield\equiv0$ \emphx{всюди в просторі}. Це \emphx{не} означає, що джерелом поля $\Dfield$ є вільні заряди --- таких зарядів у електреті просто немає! Це означає лише, що $\Dfield$ не має джерел у сенсі дивергенції --- точно так само, як завжди $\nabla\cdot\Bfield\equiv0$, хоча ніхто не стверджує, ніби через це поля $\Bfield$ не існує. Фізичною причиною самого поля $\Dfield$ (як і $\Efield$) у просторі є \emphx{зв'язані} заряди електрета; просто конкретна комбінація $\Dfield=\Efield+4\pi \vect{P}$ виявляється тотожно безджерельною за дивергенцією.
\end{Attention}

Один з поширених способів виготовлення електретів полягає в тому, що діелектрик нагрівають і одночасно поміщають у сильне електричне поле, так щоб
полярні молекули встигли переорієнтуватись і вишикуватись уздовж поля. Якщо потім діелектрик охолодити (не знімаючи поля), теплова рухливість
молекул різко падає, повернення диполів у хаотичний стан ускладнюється, і після зняття зовнішнього поля поляризація зберігається протягом місяців,
років, а іноді й десятиліть. Такий стан називають \enquote{\emphx{замороженою} поляризацією}.

Кварц та деякі інші форми діоксиду кремнію є природними електретами. У сучасній техніці більшість електретів виготовляють із синтетичних полімерів
--- фторопластів, поліпропілену, поліетилентерефталату (ПЕТ) тощо. Електрети широко застосовуються в електретних мікрофонах, дозиметрах, датчиках
тиску та інших пристроях, де потрібне джерело сталого електричного поля без витрат енергії на його підтримку.

%% --------------------------------------------------------
\section{Сегнетоелектрики}\label{sec:FE}
%% --------------------------------------------------------

\emphx{Сегнетоелектрик} --- це кристалічний діелектрик, який має \emphx{спонтанну поляризацію всередині окремих доменів}, що виникає \emphx{без}
будь-якого зовнішнього електричного поля. \emphx{Домен} --- це локальна область усередині сегнетоелектрика (типового розміру від кількох до сотень
мікрометрів), у межах якої дипольні моменти всіх елементарних комірок кристала мимовільно вишикувані в одному напрямку, тобто поляризація в домені
стала (рис.~\ref{pic:Domains}). Однак у різних доменах напрямки поляризації, як правило, різні, тому сумарна поляризація незбудженого зразка в цілому
близька до нуля.

\begin{SCfigure}[1.5][h!]\centering
	\includegraphics[width=0.4\linewidth]{domains}
	\caption{Доменна структура сегнетоелектрика\label{pic:Domains}}
\end{SCfigure}

Під впливом зовнішнього поля межі доменів (домені стінки) зміщуються, а самі домени переорієнтуються так, щоб їхня поляризація узгоджувалась із
напрямком зовнішнього поля; домени, поляризація яких збігається з полем, ростуть за рахунок сусідніх. У результаті сумарна поляризація зразка
зростає значно швидше й до значно більших значень, ніж у звичайних лінійних діелектриках.

Ця перебудова доменної структури має важливу особливість: вона \emphx{неоднозначна} і зберігає \enquote{пам'ять} про попередню історію зразка. Залежність
поляризації $P$ від поля $E$ при циклічній зміні останнього утворює характерну петлю, яку називають \emphx{петлею гістерезису}
(рис.~\ref{pic:Hysteresis}):

\begin{figure}[h!]\centering
	\localinput{hysteresis.tikz}
	\caption{Петля гістерезису сегнетоелектрика: $P_r$ --- залишкова поляризація, $E_c$ ---
		коерцитивна сила\label{pic:Hysteresis}}
\end{figure}
\begin{enumerate}
	\item за досить сильного поля $E$ поляризація виходить на насичення --- усі домени вже переорієнтовані, і подальше зростання $P$ з полем
	      відбувається практично лінійно, як у звичайного діелектрика;
	\item якщо тепер зменшувати поле до нуля, поляризація \emphx{не} повертається в нуль, а залишається на певному рівні $P_r$, який називають
	      \emphx{залишковою (остаточною) поляризацією};
	\item щоб звести поляризацію до нуля, потрібно прикласти поле протилежного знака --- $-E_c$; величину $E_c$ називають \emphx{коерцитивною
		      силою};
	\item при подальшому зростанні поля протилежного знака поляризація знову виходить на насичення у протилежний бік, а при поверненні поля до нуля
	      залишається на рівні $-P_r$, --- і цикл замикається.
\end{enumerate}
Ці властивості спостерігаються лише в певному діапазоні температур, нижчих за характерну для кожного сегнетоелектрика \emphx{температуру Кюрі}: вище
цієї температури тепловий рух руйнує спонтанну доменну структуру, і речовина переходить у звичайний (параелектричний) стан.

\begin{Attention}
	Мікроскопічне походження сегнетоелектричного стану ми вже якісно обговорили в розділі~II (п.~\ref{sec:ClausiusMossotti}): якщо добуток
	$\frac{4\pi}{3}n\alpha$ у формулі Клаузіуса--Моссотті наближається до одиниці, середовище стає нестійким щодо виникнення \emphx{спонтанної}
	поляризації --- це і є \enquote{поляризаційна катастрофа}. Не випадково сегнетоелектрики --- це, як правило, речовини з дуже великою
	поляризовністю молекул чи йонів.
\end{Attention}

Найвідоміші сегнетоелектрики --- титанат барію $\ce{BaTiO3}$ і сегнетова сіль $\ce{KNaC4H4O6*4H2O}$ (саме на честь останньої й названо весь цей клас
речовин). Завдяки аномально великим значенням діелектричної проникності $\epsilon$ (яка для сегнетоелектриків може сягати $\sim 10^3$--$10^4$)
сегнетоелектрики застосовують у мініатюрних конденсаторах великої ємності та варикондах (нелінійних конденсаторах зі змінною ємністю), а завдяки
сильному зв'язку між поляризацією й механічною деформацією --- у широкому спектрі електромеханічних і механоелектричних перетворювачів: датчиках
мікропереміщень, гідрофонах, акселерометрах, стабілізаторах частоти тощо.

%% --------------------------------------------------------
\section{П'єзоелектрики та піроелектрики}\label{sec:PiezoPyro}
%% --------------------------------------------------------

\emphx{П'єзоелектрики} --- це діелектрики (як правило, кристали без центра симетрії), в яких механічна деформація і електрична поляризація тісно
пов'язані одна з одною:
\begin{itemize}
	\item \emphx{прямий п'єзоефект} --- механічна деформація кристала (стиск, розтяг, зсув) індукує на його поверхні електричний заряд, тобто
	      створює поляризацію та різницю потенціалів;
	\item \emphx{зворотний п'єзоефект} --- прикладене зовнішнє електричне поле, навпаки, спричиняє механічну деформацію кристала.
\end{itemize}
Класичним прикладом п'єзоелектрика є кварц; на прямому й зворотному п'єзоефекті ґрунтується робота кварцових резонаторів, п'єзодатчиків тиску та
прискорення, п'єзозапальничок, ультразвукових випромінювачів і приймачів.

\begin{Attention}
	П'єзоефект, як і квадратична нелінійність $\chi^{(2)}$ (розділ~II, п.~\ref{sec:PolTensor}), можливий \emphx{лише} в кристалах без центра
	симетрії: при інверсії кристала деформація змінює знак, а поляризація --- ні, тож у центросиметричному кристалі прямий і зворотний ефекти
	взаємно компенсуються. Це ще один прояв загального правила: \enquote{нецентросиметричність --- джерело лінійних за полем ефектів}.
\end{Attention}

\emphx{Піроелектрики} --- це кристалічні діелектрики, що мають \emphx{спонтанну} поляризацію навіть за відсутності зовнішніх впливів --- подібно до
сегнетоелектричних доменів, тільки в масштабі всього кристала й без можливості переорієнтації полем. Величина цієї спонтанної поляризації
змінюється при зміні температури кристала (а також при його деформації), і ця зміна виявляється як поява (чи зміна) електричного заряду на
поверхні --- звідси й назва явища. Оскільки зміна поляризації при деформації --- це, за визначенням, п'єзоефект, \emphx{усі піроелектрики є
	п'єзоелектриками}, однак не всі п'єзоелектрики мають ще й піроелектричний ефект (для цього потрібна наявність уже готового, ненульового
спонтанного дипольного моменту).

Піроелектрики використовуються в теплових сенсорах (піроелектричних детекторах ІЧ-випромінювання, тепловізорах, датчиках руху), де зміна
температури приймача під дією випромінювання перетворюється на електричний сигнал.

%% --------------------------------------------------------
\section{Параелектрики, релаксори та антисегнетоелектрики}
%% --------------------------------------------------------

\emphx{Параелектрики.} Вище, у п.~\ref{sec:FE}, вже згадувалось, що поза межами температури Кюрі сегнетоелектрик переходить у звичайний,
\emphx{параелектричний} стан: спонтанна доменна поляризація зникає, проте поляризовність $\chi$ залишається аномально високою і швидко спадає з
температурою за законом Кюрі--Вейсса, $\chi \sim 1/(T-T_C)$. Параелектрична фаза сегнетоелектричних матеріалів (наприклад, титанату барію вище
$\sim\!120\ ^\circ\text{C}$) використовується у високоємнісних керамічних конденсаторах і термочутливих датчиках.

\emphx{Релаксори.} Окремий клас складають \emphx{релаксорні сегнетоелектрики} (релаксори) --- хімічно невпорядковані кристали (типовий приклад ---
ніобат магнію-свинцю $\ce{Pb(Mg1/3Nb2/3)O3}$, скорочено PMN), у яких замість чітких доменів утворюються нанорозмірні полярні області. Завдяки цьому
релаксори мають аномально великий і плавний (без різкого фазового переходу) відгук поляризації на поле й механічну деформацію, тому їх --- часто у
вигляді твердих розчинів PMN--PT (з титанатом свинцю) --- широко застосовують в ультразвукових перетворювачах медичної діагностики, гідроакустичних
антенах і прецизійних п'єзоприводах.

\emphx{Антисегнетоелектрики.} Наприклад, цирконат свинцю $\ce{PbZrO3}$ --- мають доменну структуру, подібну до сегнетоелектричної, проте сусідні
елементарні комірки в них поляризовані \emphx{антипаралельно}, тому за відсутності поля сумарна поляризація дорівнює нулю. Достатньо сильне
зовнішнє поле здатне індукувати перехід у сегнетоелектричний стан з великою поляризацією; після зняття поля матеріал повертається у вихідний
антисегнетоелектричний стан. Цей стрибкоподібний вихід на велику поляризацію без залишку після зняття поля робить антисегнетоелектрики зручними
для конденсаторів накопичення енергії імпульсної дії.

%% --------------------------------------------------------
\section{Рідкі кристали}\label{sec:LiquidCrystals}
%% --------------------------------------------------------

\emphx{Рідкі кристали} --- органічні речовини (як правило, з витягнутими, стрижнеподібними молекулами), що в певному інтервалі температур
утворюють проміжну (мезоморфну) фазу: зберігають плинність рідини, але їхні молекули при цьому частково впорядковані за орієнтацією --- подібно до
кристала. Через видовжену форму молекул їхня поляризовність уздовж довгої осі помітно відрізняється від поляризовності впоперек неї, тож рідкий
кристал --- наочний приклад речовини з \emphx{анізотропним} тензором поляризовності~\eqref{eq:PolTensorLinear} (розділ~II): зовнішнє поле не лише
поляризує молекули, а й \emphx{переорієнтовує} їх, розвертаючи вздовж або впоперек поля залежно від знака діелектричної анізотропії
$\chi_\parallel - \chi_\perp$. Саме на керованій полем переорієнтації молекул і пов'язаній з нею зміні оптичних властивостей середовища ґрунтується
робота рідкокристалічних дисплеїв (LCD).

%% --------------------------------------------------------
\section{Пасивні діелектрики: ізолятори та конденсаторні матеріали}
%% --------------------------------------------------------

На відміну від щойно розглянутих \emphx{активних} діелектриків (сег\-не\-то-, п'єзо-, піроелектриків, електретів), властивості яких самі по собі є
функціональними, у переважній більшості технічних застосувань діелектрик відіграє \emphx{пасивну} роль --- ізолює провідники один від одного та від
землі або слугує середовищем конденсатора. Такі \emphx{пасивні діелектрики} обирають насамперед за високою електричною міцністю (максимальним
полем, яке речовина витримує без пробою) та малими діелектричними втратами. За агрегатним станом їх поділяють на:
\begin{itemize}
	\item \emphx{тверді} --- слюда, порцеляна, скло, кераміка на основі оксиду алюмінію (ізолятори ліній електропередач, підкладки мікросхем), а
	      також полімери: поліетилен і фторопласт --- для ізоляції кабелів, полістирол і поліпропілен --- у плівкових конденсаторах,
	      епоксидно-склотканинний ламінат (FR-4) --- як основа друкованих плат;
	\item \emphx{рідкі} --- трансформаторна (мінеральна) олія, що одночасно ізолює обмотки силових трансформаторів і відводить від них тепло;
	\item \emphx{газоподібні} --- елегаз ($\ce{SF6}$), завдяки надзвичайно високій електричній міцності застосовується для ізоляції у
	      високовольтних вимикачах і закритих розподільчих пристроях. Завдяки надзвичайно високій електричній міцності чудово ізолює та гасить електричну дугу.
\end{itemize}

%%% --------------------------------------------------------
%\section{Порівняльна таблиця типів діелектриків}
%%% --------------------------------------------------------
%
%Щоб систематизувати все сказане вище, зведемо основні типи діелектриків у таблицю~\ref{tab:DielectricTypes}.
%
%\begin{table}[h!]
%	\centering
%	\small
%	\caption{Порівняння типів діелектриків\label{tab:DielectricTypes}}
%\DefTblrTemplate{contfoot-text}{default}{}
%\DefTblrTemplate{conthead-text}{default}{}
%\DefTblrTemplate{caption}{default}{}
%\DefTblrTemplate{conthead}{default}{}
%\DefTblrTemplate{capcont}{default}{}
%\begin{longtblr}[]{
%	colspec = {Q[c,m,2.4cm] c c c Q[c,m,3.4cm]},
%    rowhead=1,
%	row{1} = {bg=rowfill, fg=themecolordark,  c},
%	hlines = {1-Z}{infotablecolor},
%	vlines = {1-Z}{infotablecolor},
%	}
%		Тип & {Спонтанна \\ поляризація?} & {Керованість \\ полем?} & $\epsilon$ & Застосування \\
%		Лінійний діелектрик   & ні              & {так\\ (пропорційно)}   & $2$--$10$         & ізолятори, конденсатори \\
%		Електрет               & {так\\ (заморожена)} & ні                 & ---               & мікрофони, дозиметри \\
%		Сегнетоелектрик        & {так\\ (домени)}     & {так\\ (гістерезис)}   & $10^3$--$10^4$    & вариконди, датчики \\
%		П'єзоелектрик          & ні               & {так\\ (деформація)}   & $5$--$100$        & резонатори, датчики \\
%		Піроелектрик           & так              & {ні\\ (тільки $T$)}    & ---               & тепловізори, ІЧ-сенсори \\
%		Релаксор               & {так\\ (нанообласті)} & так (плавно)      & $10^3$--$10^4$    & ультразвук, приводи \\
%		Антисегнетоелектрик    & {так\\ (антипарал.)} & {так\\ (стрибком)}     & велика            & накопичувачі енергії \\
%		Рідкий кристал         & ні               & {так\\ (переорієнтація)} & анізотропна    & LCD-дисплеї \\
%	\end{longtblr}
%\end{table}
%
%\begin{Attention}
%	Зверніть увагу на спільну рису всіх \enquote{активних} діелектриків: їхні властивості --- спонтанна поляризація, гістерезис, п'єзоефект,
%	піроефект --- виникають \emphx{без} зовнішнього поля або \emphx{нелінійно} залежать від нього. Це принципово відрізняє їх від звичайних
%	лінійних діелектриків, для яких $\vect P = \chi\Efield$ і $\chi$ --- стала. Саме ця нелінійність і робить їх функціональними матеріалами
%	сучасної техніки.
%\end{Attention}